\chapter{Model-based reinforcement learning and planning}
\label{ch:model-based}
\label{ch:search}

Model-free methods learn values or policies directly from experience.
Model-based methods also use a \emph{model} of the environment, given or
learnt, to \emph{plan}: to compute what to do by simulating rather than by
acting. Planning can happen in the background, generating simulated
experience to train the values or the policy (Dyna, MBPO), or at decision
time, searching from the current state before each move (rollouts, MCTS,
AlphaZero, MuZero). This chapter covers both, then returns to Pacman: the
look-ahead search of the repository, the question of what the search may
know about the ghosts, and expert iteration, which teaches the search's
decisions back to the network.

\section{Models and background planning}

A \emph{model} predicts the next state and reward, $\hat p(s', r \mid s,
a)$, either as a distribution or as a sampler. It can be known (the rules
of a board game, a physics simulator, a copy of the Pacman environment) or
learnt from the transitions observed, by supervised learning.

\begin{algobox}{Dyna-Q (Sutton, 1990)}
  \begin{algorithmic}[1]
    \For{each real step}
      \State in state $s$, play $a$ ($\varepsilon$-greedy w.r.t.\ $Q$),
        observe $r, s'$
      \State Q-learning update with $(s, a, r, s')$
      \State model update: $\mathrm{Model}(s, a) \gets (r, s')$
      \For{$n$ planning steps}
        \State pick a previously visited $(\tilde s, \tilde a)$;
          $(\tilde r, \tilde s') \gets \mathrm{Model}(\tilde s, \tilde a)$
        \State Q-learning update with $(\tilde s, \tilde a, \tilde r,
          \tilde s')$
      \EndFor
    \EndFor
  \end{algorithmic}
\end{algobox}

Each real transition is reused through the model many times, propagating
values much faster than with Q-learning alone. \emph{Prioritised sweeping}
chooses the simulated updates where they matter: when a value changes, the
predecessors of the state are queued with a priority equal to the size of
the change they would undergo.

A learnt model is wrong, and planning amplifies its errors: the planner
seeks the actions whose predicted outcome is best, including those where
the model is optimistically wrong. The following lemma quantifies the
effect of model errors on values.

\begin{lemma}{Simulation lemma}{simulation}
  Let $M = (p, r)$ and $\hat M = (\hat p, \hat r)$ be two MDPs on the same
  states and actions with rewards bounded by $R_{\max}$, and suppose
  $\abs{r(s, a) - \hat r(s, a)} \leq \varepsilon_r$ and $\sum_{s'}
  \abs{p(s' \mid s, a) - \hat p(s' \mid s, a)} \leq \varepsilon_p$ for all
  $s, a$. Then for every policy $\pi$,
  \begin{equation}
    \norm{v^\pi_M - v^\pi_{\hat M}}_\infty \leq \frac{\varepsilon_r}{1 -
    \gamma} + \frac{\gamma\, \varepsilon_p\, R_{\max}}{(1 - \gamma)^2}.
  \end{equation}
\end{lemma}

\begin{proof}
  With the matrix notation of \cref{prop:matrix}, $v = r_\pi + \gamma P_\pi
  v$ and $\hat v = \hat r_\pi + \gamma \hat P_\pi \hat v$, so
  \begin{equation*}
    v - \hat v = (r_\pi - \hat r_\pi) + \gamma (P_\pi - \hat P_\pi) v +
    \gamma \hat P_\pi (v - \hat v).
  \end{equation*}
  Each row of $P_\pi - \hat P_\pi$ has $\ell_1$ norm at most
  $\varepsilon_p$ and $\norm{v}_\infty \leq R_{\max} / (1 - \gamma)$
  (\cref{prop:return-props}), and $\hat P_\pi$ is stochastic. Taking sup
  norms, $\norm{v - \hat v}_\infty \leq \varepsilon_r + \gamma
  \varepsilon_p R_{\max} / (1 - \gamma) + \gamma \norm{v - \hat
  v}_\infty$; solve.
\end{proof}

\begin{intuition}
  A model that is wrong about the rewards by $\varepsilon_r$ makes that
  error at every simulated step, hence $\varepsilon_r / (1 - \gamma)$ over
  the horizon. A model that is wrong about the transitions is worse: at
  each step it sends the simulation to slightly wrong states, and landing
  in the wrong place can cost as much as the whole range of values,
  $R_{\max} / (1 - \gamma)$; this too happens at every step, hence the
  square. Imagined trajectories drift away from reality the longer they
  run.
\end{intuition}

The factor $1 / (1 - \gamma)^2$ reflects the compounding of errors over
long simulated trajectories. Modern methods therefore keep the simulated
trajectories short: MBPO (Janner et al., 2019) branches rollouts of a few
steps from real states with an ensemble of learnt models, and trains SAC
on the mixture of real and simulated transitions. World-model methods
(Dreamer, Hafner et al., 2020--2023) learn a latent dynamics model from
images and train an actor-critic entirely in imagination.

\section{Decision-time planning}

Background planning improves the policy everywhere; decision-time planning
concentrates the computation on the current state. It is how game-playing
programs have always worked, and it pays off when the policy alone is
imperfect: a trained policy decides each move in one glance at the board,
and when the PPO agent of \cref{ch:ppo} loses, it is usually because it
walks into a trap that only closes a few moves later.

\subsection{Rollouts}

Suppose we can simulate the game from the current state: here, a copy of
the environment (\code{PacmanEnvironment.clone()}). To evaluate a move $a$
in state $s$, play $a$ in a copy, then let a \emph{rollout policy} $\pi$
play for at most $H$ moves, and score the simulation by
\begin{equation}
  \label{eq:rollout-return}
  \hat G = \sum_{k=0}^{m-1} \gamma^k R_{k+1} + \gamma^m\, V(S_m),
\end{equation}
where $m \leq H + 1$ is the number of moves simulated and $V(S_m) = 0$ if
the game ended. Averaging $N$ independent simulations gives an estimate
$\hat q(s, a)$.

\begin{proposition}{What a rollout estimates}{rollout}
  If $V = v_\pi$, then $\E[\hat G] = q_\pi(s, a)$ for every horizon $H$,
  and the average of $N$ independent simulations has variance
  $\mathrm{Var}(\hat G) / N$. Playing in $s$ an action maximising $q_\pi(s,
  \cdot)$, and following $\pi$ afterwards, is at least as good as $\pi$.
\end{proposition}

\begin{proof}
  $\E[\hat G]$ is $(\T^\pi)^{H} V$ applied after the first move, as in
  \cref{prop:nstep-error}; with $V = v_\pi$ every application of $\T^\pi$
  leaves $v_\pi$ unchanged, so $\E[\hat G] = \E[R_1 + \gamma v_\pi(S_1)
  \mid s, a] = q_\pi(s, a)$. The variance of a mean of $N$ independent
  copies is the variance divided by $N$. The last claim is the policy
  improvement theorem (\cref{cor:greedy-improvement}) applied in the single
  state $s$.
\end{proof}

In the repository the rollout policy is the trained agent itself (greedy
moves), and $V$ its value head; the rollouts therefore estimate what
happens if the agent plays $a$ and then goes on as usual.

\paragraph{Common random numbers.} The moves of the ghosts are random. To
compare two candidate moves, the $k$-th simulation of each move uses the
same random numbers. For two estimates $X$ and $Y$,
\begin{equation}
  \mathrm{Var}(X - Y) = \mathrm{Var}(X) + \mathrm{Var}(Y) - 2\,
  \mathrm{Cov}(X, Y),
\end{equation}
and simulations that share the ghosts' random draws are positively
correlated, so their \emph{difference}, which decides the move, is much
less noisy than each of them.

\subsection{Two ways of choosing the move}

\begin{definition}{Value mode and safety mode}{search-modes}
  Let $M$ be the set of candidate moves (the moves leading to different
  cells), $a_0$ the agent's move, and, for each $a \in M$, $\hat q(a)$ the
  average score and $\hat d(a)$ the fraction of its $N$ simulations in
  which Pacman dies.
  \begin{itemize}[nosep]
    \item \emph{Value mode:} play $\argmax_{a \in M} \hat q(a)$.
    \item \emph{Safety mode:} keep $a_0$ unless $\hat d(a_0) - \min_a \hat
      d(a) \geq \kappa$; otherwise play, among the moves of minimal death
      rate, the one of best average score. The repository uses $\kappa =
      0.15$.
  \end{itemize}
  The search is only used when a ghost is at most $5$ cells away
  (Manhattan distance), about $33$ moves per game.
\end{definition}

\begin{pacman}{Why the value mode fails}{value-mode}
  With $N = 20$ simulations, the value mode \emph{lowered} the win rate
  (about $95.5\%$, against $96.1\%$ without search) and changed about a
  quarter of the agent's moves. The reason is noise: the returns of a
  simulation range from $-35$ (death) to $+20$ and more (a ghost eaten),
  so the standard deviation of a mean of $20$ of them is of the order of
  a few points, while the difference in value between two reasonable moves
  is often a fraction of a point. The maximum of noisy estimates then picks
  a random move (\cref{prop:max-bias}). The death rate is a much more
  robust signal: a trap shows up as many deaths for one move and few for
  another. The safety mode changes the agent's move about $0.6$ times per
  game, only when it clearly walks into a trap.
\end{pacman}

\begin{intuition}
  The safety mode does not ask the search which move is best, a question
  its $20$ noisy simulations per move cannot answer reliably. It asks a
  simpler one: is the move the agent wants to play clearly more dangerous
  than another? The trained policy keeps the initiative, and the search
  acts as a safety net that vetoes a move only when the evidence is strong.
\end{intuition}

\begin{lstlisting}[caption={Choosing the move (\code{pacman\_search.py},
  \code{LookaheadSearch.select\_action}, excerpt)}]
scores = {m: np.mean([s["score"] for s in sims if s["move"] == m])
          for m in moves}
if self.mode == "value":
    best = max(moves, key=lambda m: scores[m])
    if scores[agent_move] >= scores[best] - 1e-9:
        return agent_action
else:  # safety mode
    deaths = {m: np.mean([s["env"].dead for s in sims
                          if s["move"] == m]) for m in moves}
    safest = min(deaths.values())
    if deaths[agent_move] - safest < self.margin:   # kappa = 0.15
        return agent_action
    best = max((m for m in moves if deaths[m] == safest),
               key=lambda m: scores[m])
return best
\end{lstlisting}

\section{Monte Carlo tree search}

The search of the repository is the simplest form of Monte Carlo planning:
a tree of depth one with the same number of simulations for every move.
Monte Carlo tree search (MCTS) grows a deeper tree, asymmetrically, and
allocates the simulations to the moves that deserve them. Each node of the
tree stores, for each action, a visit count $N(s, a)$ and a mean value
$\hat q(s, a)$.

\begin{algobox}{Monte Carlo tree search (one simulation)}
  \begin{algorithmic}[1]
    \State \textbf{selection:} from the root, descend the tree choosing at
      each node the action maximising a selection rule (UCT below), until a
      leaf
    \State \textbf{expansion:} add one child of the leaf to the tree
    \State \textbf{evaluation:} estimate the value of the new node, by a
      rollout with a default policy, or by a value network
    \State \textbf{backup:} along the path, increment $N(s, a)$ and update
      $\hat q(s, a)$ with the observed return
  \end{algorithmic}
\end{algobox}

After the simulations, the move played is the most visited action of the
root. The UCT selection rule (Kocsis and Szepesv\'ari, 2006) treats each
node as a bandit and applies UCB1 (\cref{def:ucb1}):
\begin{equation}
  a = \argmax_{a} \left( \hat q(s, a) + c \sqrt{\frac{\ln N(s)}{N(s,
  a)}} \right),
\end{equation}
with $N(s) = \sum_a N(s, a)$. Kocsis and Szepesv\'ari proved that, with
suitable $c$, the probability of choosing a wrong move at the root tends
to $0$ as the number of simulations grows. In a stochastic game such as
Pacman, the tree is built over sequences of the agent's moves
(``open-loop'' MCTS), each simulation replaying the sequence on a fresh
copy of the game, or the tree also branches on chance outcomes.

\begin{intuition}
  MCTS places a bandit at every node of the tree. Moves that look good are
  simulated more and their subtrees grow deeper; moves that look bad are
  revisited only now and then, in case their first results were unlucky.
  The tree grows towards the promising lines, as a chess player analyses
  the critical variations deeply and dismisses the others after a glance.
\end{intuition}

\begin{remark}{MCTS for two-player games}{mcts-two-players}
  In a game between two opponents, such as Hex, chess or Go, the players
  alternate along each path of the tree, and the statistics of a node are
  kept from the point of view of the player who made the move leading to
  it. During the backup the result is therefore flipped at each level: a
  win for one player is a loss for the other. Each player then selects,
  with UCT, the moves that are best \emph{for itself}, which is the minimax
  principle obtained by sampling. The move finally played can be the most
  visited child of the root (the usual, most robust choice), the child with
  the highest mean value (optimistic when it has few visits), or the child
  with the highest lower confidence bound (pessimistic).

  A C++ implementation for Hex, with brute force, flat Monte Carlo and MCTS
  players, and a step-by-step account of the algorithm on a small example,
  is in the repository
  \href{https://github.com/CL-BZH/Hex_game_AI}{CL-BZH/Hex\_game\_AI}
  (\code{hex\_machine\_mcts.h} and \code{Docs/mcts.pdf}). The
  \href{https://CL-BZH.github.io/MCTS-Visualizer/}{MCTS Visualizer}
  (repository
  \href{https://github.com/CL-BZH/MCTS-Visualizer}{CL-BZH/MCTS-Visualizer})
  shows the tree growing iteration by iteration on a simple two-player game,
  with the four phases and the UCB1 values of each node.
\end{remark}

\subsection{AlphaZero}

AlphaZero (Silver et al., 2018) learnt chess, shogi and Go from the rules
alone by combining MCTS with a network $f_\theta(s) = (\mathbf{p}, v)$ that
outputs a prior policy and a value. MCTS uses them through the PUCT rule
\begin{equation}
  a = \argmax_a \left( \hat q(s, a) + c\, p_a(s) \frac{\sqrt{N(s)}}{1 +
  N(s, a)} \right),
\end{equation}
which explores first the moves the network finds promising, and evaluates
the leaves with $v$ instead of rollouts. The games of self-play provide the
training data: the network is trained to predict the visit distribution of
the search, $\boldsymbol{\pi}_t \propto N(s_t, \cdot)$, and the final
result $z$ of the game,
\begin{equation}
  \ell(\theta) = (z - v)^2 - \boldsymbol{\pi}_t^\top \log \mathbf{p} + c
  \norm{\theta}^2.
\end{equation}
The search is a \emph{policy improvement operator}: its visit distribution
is better than the prior it started from, and training the network on it
is \emph{expert iteration} (Anthony et al., 2017). The network then makes
the next searches better, and so on.

\begin{intuition}
  In PUCT the prior $p_a(s)$ plays the role of intuition: the moves the
  network likes are explored first, and the bonus, decreasing like $1 /
  (1 + N(s, a))$, lets the search's own statistics take over as they
  accumulate. In the loss, the value head learns to predict who wins, and
  the policy head learns to imitate the search, which plays better than
  the raw network because it looks ahead. The program pulls itself up: a
  better network makes a better search, which makes a better teacher.
\end{intuition}

\subsection{MuZero}

MuZero (Schrittwieser et al., 2020) plans without knowing the rules. It
learns three functions: a representation $h$ mapping observations to a
latent state, a dynamics function $g(z, a) = (\hat r, z')$, and a
prediction function $f(z) = (\mathbf{p}, v)$. The latent model is not trained
to reconstruct observations, only to predict, along real trajectories
unrolled for a few steps, the rewards, the values and the search policies:
it models only what matters for planning. MuZero matched AlphaZero in board
games and set the state of the art on Atari at the time.

\section{What does the search know about the ghosts?}

A copy of the environment simulates the ghosts with the game's own code,
which the agent never sees. Using it would give the search a \emph{perfect
model} of the opponents: legitimate in a board game whose rules are known,
debatable for Pacman, where the behaviour of the ghosts is part of what the
agent should discover. The search therefore offers three \emph{ghost
models}:
\begin{itemize}[nosep]
  \item \code{exact}: the game's code (perfect model), used only as a
    reference;
  \item \code{chase} (default): worst case; when Pacman has no power each
    ghost moves along a shortest path towards him, and randomly when he has
    power;
  \item \code{estimated}: a ghost chases with probability $p_{\mathrm{chase}}$
    and moves randomly otherwise; when Pacman has power it moves away from
    him with probability $p_{\mathrm{flee}}$. Both probabilities are
    estimated for each ghost from its observed moves.
\end{itemize}
With \code{chase} and \code{estimated}, the search uses only what the agent
observes (positions, power) and the rules for Pacman himself.

\begin{proposition}{Estimating a ghost's chase probability}{ghost-estimator}
  Model each observed move of a ghost as: towards Pacman with probability
  $p$, otherwise uniform among its $n$ reachable cells, of which $m$ go
  towards Pacman ($0 < m < n$). Then the probability that the move goes
  towards Pacman is $p + (1 - p)\, m / n$, and after observing moves with
  ratios $q_i = m_i / n_i$ the posterior density of $p$ under a prior
  $\mathrm{Beta}(2, 2)$ is proportional to
  \begin{equation}
    p (1 - p) \prod_{i \in \mathcal{T}} \left( p + (1 - p) q_i \right)
    \prod_{i \notin \mathcal{T}} (1 - p)(1 - q_i),
  \end{equation}
  where $\mathcal{T}$ is the set of moves that went towards Pacman.
\end{proposition}

\begin{proof}
  A move goes towards Pacman either because the ghost chases (probability
  $p$) or because it moves randomly (probability $1 - p$) and picks one of
  the $m$ cells out of $n$. The moves are independent given $p$, and the
  posterior is the prior density $\propto p(1 - p)$ times the likelihood.
\end{proof}

\begin{intuition}
  A move towards Pacman is evidence that the ghost chases, but weak
  evidence when most of the ghost's possible moves go towards Pacman anyway
  ($q_i$ close to $1$): a random ghost would have done the same. A move
  away from Pacman is strong evidence of randomness, since a chasing ghost
  never makes it. The posterior weighs each observation by how much it
  discriminates between the two behaviours.
\end{intuition}

The repository takes the maximum of this posterior on a grid. Over many
games, it recovers $p \approx 0.73$ and $0.64$ for the two ghosts (their
configured aggressiveness is $0.74$ and $0.64$), and a flee probability of
about $0.08$: edible ghosts move almost at random. The first version of
the models assumed that edible ghosts always flee; the simulated hunts then
failed, the search judged eating a ghost dangerous and made the agent
hesitate in front of an edible ghost. Estimating the flee probability
fixed it: a model that is wrong in the optimistic or the pessimistic
direction can both mislead a planner.

\begin{pacman}{The search on 1000 games}{search-results}
  With $20$ simulations, horizon $15$, safety mode, the search was
  evaluated on $1000$ games: all the games lost without search, and a
  sample of the won ones, were replayed with search from the same starting
  conditions.
  \begin{center}
    \begin{tabular}{@{}lc@{}}
      \toprule
      Ghost model & Estimated win rate \\
      \midrule
      no search & $96.1\%$ \\
      \code{exact} & $97.1\% \pm 0.5$ \\
      \code{chase} & $97.5\% \pm 0.5$ \\
      \code{estimated} & $97.7\% \pm 0.3$ \\
      \bottomrule
    \end{tabular}
  \end{center}
  The models that do not know the ghosts' code do at least as well as the
  exact one: the gain comes from seeing a trap a few moves ahead, not from
  knowing the opponents' rules.
\end{pacman}

\section{Expert iteration}

The search is expensive (up to about $100$ simulated games per searched
move). If
the network could learn to make the search's decisions by itself, the gain
would be free at play time. Simply training the network on games played
\emph{by the search} (behaviour cloning) has a known flaw: the network only
sees the situations that the search gets into, and when it makes a mistake
of its own, it reaches situations it has never seen, where its errors
compound (Ross and Bagnell, 2010). The remedy, DAgger (Ross, Gordon and
Bagnell, 2011), lets the \emph{learner} play and asks the \emph{expert}
what it would have done in the situations the learner reached.

\begin{algobox}{Expert iteration (DAgger form), \code{pacman\_expert.py}}
  \begin{algorithmic}[1]
    \State start from the PPO agent $\pi_\theta$
    \For{each iteration}
      \State the agent plays $2000$ games alone; whenever a ghost is close,
        record the state $s$ and the search's move $a^\star(s)$
      \State record also a fraction of the other states, labelled with the
        agent's own move, and the discounted returns observed
      \State train on the data of the last iterations: weighted
        cross-entropy $-w(s) \log \pi_\theta(a^\star(s) \mid s)$ with
        $w(s) = 10$ where the search disagreed with the agent and $1$
        otherwise, and squared error of the value head
    \EndFor
  \end{algorithmic}
\end{algobox}

The labels of the states where the search agrees with the agent keep the
rest of its behaviour unchanged; the weight $10$ makes the few
disagreements count. Because the games are played by the agent alone,
each iteration also measures its progress.

\begin{pacman}{Results of expert iteration}{expert-results}
  Ten iterations of $2000$ games on Kaggle (search with $10$ simulations,
  horizon $10$, \code{chase} model), about one hour each, almost all of it
  spent in the search (on CPU). The search disagreed with the agent
  $0.27$ times per game at the first iteration and $0.12$ at the last, and
  the network reproduced $76\%$ of the corrections of its training data.
  On the same $1000$ games:
  \begin{center}
    \begin{tabular}{@{}lcc@{}}
      \toprule
      Same $1000$ games & Without search & With search ($20$) \\
      \midrule
      PPO agent (\code{PPO\_1}) & $97.1\%$ & $98.1\%$ \\
      After expert iteration (\code{Expert\_1}) & $97.7\%$ & $97.9\%$ \\
      \bottomrule
    \end{tabular}
  \end{center}
  Expert iteration and the search correct the same mistakes: without
  search the expert agent recovers about $60\%$ of the search's gain (and
  no longer loses games by the ``off-by-one'' of \cref{sec:game}), but with
  the search both agents reach about $98\%$.
\end{pacman}

\section{What is left?}

The repository classifies the lost games by cause (\code{pacman\_play.py}).
For the PPO agent alone, out of about $38$ losses in $1000$ games, $2$ come
from the off-by-one, $1$ from a ghost catching Pacman just after his power
ended, $6$ from \emph{bad starts} (Pacman caught in the first five moves),
and the others from traps. With the search, the same three bad starts are
lost in every configuration, which suggests that they cannot be won. The
$15$ to $18$ remaining losses resist both the search and expert iteration:
they are probably traps from which the agent's own policy, which plays the
simulations, cannot escape either. Going further would require a deeper
search (MCTS), or, following AlphaZero, alternating the improvement of the
search and of the network over many more iterations.

\section{Summary}

\begin{itemize}[nosep]
  \item A model can be used to generate simulated experience (Dyna, MBPO,
    Dreamer) or to search at decision time (rollouts, MCTS); the errors of
    a learnt model compound with the length of the simulations.
  \item MCTS treats each node as a bandit (UCT); AlphaZero guides it with a
    policy and value network and trains the network on the search's
    decisions; MuZero learns a latent model for planning.
  \item On Pacman, a depth-one search with a safety rule raises the win
    rate from $96\%$ to about $98\%$, without knowing the ghosts' code;
    expert iteration transfers part of this gain to the network alone.
\end{itemize}
